\subsection{COMPONENTS}

Given a point x of a topological space X, the union of the connected subsets of X which contain x is connected (no. 1, Proposition 2); it is therefore the largest connected subset of X which contains x.

DEFINITION 3. The component (or connected component) of a point of a topological space X is the largest connected subset of X which contains this point. The components of a subset A of X are the components of the points of A, relative to the subspace A of X.

If a space is connected, the component of each point is the whole space. If a space X is such that for each pair \((x, y)\) of points of X there is a connected set containing x and y, then X is connected.

A space X is said to be totally disconnected if the component of each point of X consists of the point alone. A subset A of X is a totally disconnected set if the subspace A of X is totally disconnected.

A discrete space is totally disconnected, but one should beware of confusing these two notions; for example, we shall see in Chapter IV, § 2, no. 5, that the rational line, which is not a discrete space, is totally disconnected.

\pandocbounded{\includegraphics[keepaspectratio]{images/part001_10033d49.jpg}}

A set which is both open and closed contains the component of each of its points, whence the component of a point is contained in the intersection of the sets which are both open and closed and which contain this point. However, the component of a point is not necessarily equal to this intersection (cf.~Exercise 9 and Chapter II, § 4, no. 4, Proposition 6).

PROPOSITION 9. The component of any point in a topological space X is a closed set. The relation ``y belongs to the component of x'' is an equivalence relation R\{x, y\} on X, and the equivalence classes are the components of X. The quotient space X/R is totally disconnected.

The first part of the proposition is an immediate consequence of Definition 3 and the fact that the closure of a connected set is connected (no. 1, Proposition 1). Since the union of connected sets which have a point in common is connected (no. 1, Proposition 2), the relation R is transitive, hence is an equivalence relation (since it is obviously reflexive and symmetric) and the equivalence class of x with respect to R is the component of x. It remains to show that X/R is totally disconnected. Let \(f: X \to X/R\) be the canonical mapping, and let F be a closed set in X/R containing at least two distinct points; the inverse image \(\overline{f}^{1}(F)\) of F is closed in X, saturated with respect to R, and contains at least two distinct components of X and hence is not connected. Hence there exist two non-empty closed sets B, C in X such that \(B \cap C = \varnothing\) and \(B \cup C = \overline{f}^{1}(F)\) . The component of any point x of \(\overline{f}^{1}(F)\) in \(\overline{f}^{1}(F)\) is the same as the component of x in X (by the definition of R) and therefore B and C, which are both open and closed in \(\overline{f}^{1}(F)\) , are saturated with respect to R. Hence \(f(B)\) and \(f(C)\) are closed in X/R, and \(f(B) \cup f(C) = F\) and \(f(B) \cap f(C) = \varnothing\) ; this shows that F is not connected and consequently that X/R is totally disconnected.

PROPOSITION 10. In a product space \(\mathbf{X} = \prod_{i\in J}\mathbf{X}_i\) , the component of \(x = (x_i)\) in \(\mathbf{X}\) is the product of the components of \(x_i\) in the factors \(\mathbf{X}_i\) .

This product set is connected (no. 4, Proposition 8). Conversely, if A is a connected subset of X which contains x, then \(\mathrm{pr}_{i}(A)\) is a connected set (no. 2, Proposition 4) which contains x; since \(A \subset \prod_{i} \mathrm{pr}_{i}(A)\) , it follows that A is contained in the product of the components of the \(x_{i}\) .

